\documentclass[11pt,a4paper]{article}
\usepackage[T1]{fontenc}
\usepackage{lmodern,amsmath,amssymb,amsthm,booktabs}
\usepackage[margin=2.5cm]{geometry}
\usepackage[hidelinks]{hyperref}

\newtheorem{theorem}{Theorem}
\newtheorem{lemma}[theorem]{Lemma}
\newtheorem{corollary}[theorem]{Corollary}
\theoremstyle{remark}
\newtheorem{remark}[theorem]{Remark}

\newcommand{\ZZ}{\mathbb{Z}}
\newcommand{\Zp}{\mathbb{Z}_{271}}

\title{A projective plane of order 28 admits no collineation of order 271}
\author{Yudai Kobayashi}
\date{November 2026\\[4pt]\small Computer-assisted; not peer reviewed.}

\begin{document}
\maketitle

\begin{abstract}
We show by an exhaustive computer search that a projective plane of order $28$ has no collineation of order $271$. Such a collineation would give a planar difference family in $\Zp$ with blocks of sizes $13$, $9$, $7$, and there is none. Together with the norm-parity lemma and known results on cyclic planes, it follows that for every order $n<36$ that is not a prime power, a projective plane of order $n$ has no collineation of prime order $p>n+1$.
\end{abstract}

\section{Introduction}

The existence of a projective plane of order $28$ is open; $28$ is not excluded by the Bruck--Ryser theorem~\cite{BR}. Since $28^2+28+1=813=3\cdot271$, the prime $271$ is the only prime $p>29$ that can be the order of a collineation (Lemma~\ref{l:fpf}).

\begin{theorem}\label{t:main}
A projective plane of order $28$ has no collineation of order $271$.
\end{theorem}

\begin{corollary}\label{c:small}
Let $n<36$ not be a prime power, and let $p>n+1$ be a prime. Then a projective plane of order $n$ has no collineation of order $p$.
\end{corollary}

For $n=36$ the prime $p=43$ divides $36^2+36+1=1333=31\cdot43$, and we do not know how to treat it.

\emph{Status.} The proof of Theorem~\ref{t:main} rests on a computer search (Section~\ref{s:search}). It was repeated with a second, separately written program, which gives the same counts at every stage. Both programs were written with the same normalisation and by the same means, so they are not fully independent. Most derivations and programs were made by AI (Claude, Anthropic) under the author's direction. We found no statement of Theorem~\ref{t:main} or Corollary~\ref{c:small} in the literature we could access, and we make no claim that they are new.

\section{Reduction to a difference family}

Let $\Pi$ be a projective plane of order $n$ with $v=n^2+n+1$ points, and $\sigma$ a collineation of prime order $p$.

\begin{lemma}\label{l:fpf}
If $p>n+1$, then $\sigma$ fixes no point and no line, and $p\mid v$.
\end{lemma}
\begin{proof}
Let $P$ be a fixed point. The $n+1<p$ lines through $P$ are permuted by $\sigma$ in orbits of length $1$ or $p$, hence are all fixed. Each of them has $n+1<p$ points, which are therefore fixed. These lines cover $\Pi$, so $\sigma=1$. Dually, no line is fixed. All point orbits have length $p$, so $p\mid v$.
\end{proof}

Now let $n=28$ and $p=271$. There are $R=v/p=3$ point orbits and $3$ line orbits. Choose representatives $x_1,x_2,x_3$ and $y_1,y_2,y_3$ and put
\[
D_{ab}=\{\,i\in\Zp : \sigma^i x_a\in y_b\,\}.
\]
Then $\sigma^i x_a\in\sigma^k y_b$ if and only if $i-k\in D_{ab}$.

\begin{lemma}\label{l:fam}
The sets $D_{11},D_{12},D_{13}$ have sizes $13,9,7$ in some order, and every nonzero element of $\Zp$ can be written as $x-y$ with $x\ne y$ in a common $D_{1b}$ in exactly one way. We call such a triple a \emph{planar difference family}.
\end{lemma}
\begin{proof}
Let $i\ne j$. The points $\sigma^ix_1$, $\sigma^jx_1$ lie on exactly one line $\sigma^ky_b$, that is, there is exactly one pair $(b,k)$ with $i-k,\,j-k\in D_{1b}$. Hence $i-j$ has exactly one representation as a difference within a block. The line $\sigma^ky_b$ passes through $x_1$ if and only if $-k\in D_{1b}$, so $\sum_b|D_{1b}|=n+1=29$, and counting differences gives $\sum_b|D_{1b}|(|D_{1b}|-1)=270$. The only solution in non-negative integers is $\{13,9,7\}$.
\end{proof}

So Theorem~\ref{t:main} follows from:

\begin{theorem}\label{t:fam}
There is no planar difference family in $\Zp$ with blocks of sizes $13$, $9$, $7$.
\end{theorem}

\section{The search}\label{s:search}

Let $\{A,B,C\}$ be a planar difference family with $|A|=13$, $|B|=9$, $|C|=7$. For $u\in\Zp^*$ and $t_A,t_B,t_C\in\Zp$, the family $\{uA+t_A,\,uB+t_B,\,uC+t_C\}$ is again a planar difference family. We use this freedom as follows.

For distinct $x,y,z\in A$ put $\lambda=(z-x)/(y-x)\notin\{0,1\}$. Reordering the triple replaces $\lambda$ by an element of its class $\{\lambda,1/\lambda,1-\lambda,1/(1-\lambda),\lambda/(\lambda-1),(\lambda-1)/\lambda\}$. The class of $-1$ is $\{-1,2,1/2\}$; a triple of this class is a $3$-term arithmetic progression $\{w-\delta,w,w+\delta\}$, which represents $\delta$ twice and cannot occur in $A$. The other classes are $45$ in number ($44$ of size $6$ and one of size $2$, since $271\equiv1\pmod 3$). We list them in a fixed order $\Lambda_1,\dots,\Lambda_{45}$ and fix $\lambda_s\in\Lambda_s$.

\begin{lemma}\label{l:norm}
Up to the operations above we may assume:
\begin{enumerate}
\item for some $s$, $\{0,1,\lambda_s\}\subseteq A$ and $A$ contains no triple of class $\Lambda_1,\dots,\Lambda_{s-1}$;
\item let $d$ be the least element of $\{1,\dots,270\}$ that is not a difference of $A$; the block among $B,C$ that has $d$ as a difference contains $\{0,d\}$;
\item let $e$ be the least element not a difference of $A$ or of that block; the remaining block contains $\{0,e\}$.
\end{enumerate}
\end{lemma}
\begin{proof}
(1) Let $s$ be least such that $A$ has a triple of class $\Lambda_s$, and order it as $(x,y,z)$ with $(z-x)/(y-x)=\lambda_s$. Apply $u=1/(y-x)$ and $t_A=-x/(y-x)$. (2), (3) Translate $B$ and $C$; this does not change the differences. The element $d$ exists because $A$ has only $156<270$ differences.
\end{proof}

The program \texttt{search.cpp} enumerates, for each $s$, all $13$-sets $A\supseteq\{0,1,\lambda_s\}$ with distinct differences, no $3$-term arithmetic progression and no triple of class $\Lambda_1,\dots,\Lambda_{s-1}$, and for each such $A$ all completions $B,C$ normalised as in Lemma~\ref{l:norm}. A candidate element is discarded only if it would repeat a difference, create an arithmetic progression or a triple of an earlier class, or if too few candidates remain. The work is split into $40$ chunks by the first two elements of $A$ outside $\{0,1,\lambda_s\}$.

\paragraph{Result.} All $40$ chunks finished. The program examined $6\,871\,065\,175$ admissible sets $A$ and found no family. The run took $11.0$ hours on $4$ threads (g++ 16.1, MinGW-w64). The per-chunk counts are in \texttt{code/results.txt}. This proves Theorem~\ref{t:fam}, and hence Theorem~\ref{t:main}, assuming the program is correct.

\paragraph{Re-check.} A second program, \texttt{recheck.cpp}, implements the same normalisation without sharing code with \texttt{search.cpp} (its own computation of the classes $\Lambda_s$, its own pruning and completion, a different splitting of the work). It also finds no family, and the number of admissible sets $A$ agrees with \texttt{search.cpp} in each of the $45$ stages (\texttt{code/stage\_counts.txt}).

\paragraph{Controls.} The same program, with parameters for $n=13$, $p=61$ (blocks $7,4,3$) and $n=19$, $p=127$ (blocks $9,7,4$), finds $28$ and $8$ normalised families, which fall into $12$ and $4$ classes under the operations above. Both cases occur in the Desarguesian planes $\mathrm{PG}(2,13)$ and $\mathrm{PG}(2,19)$, where the Singer cycle has a subgroup of order $61$, respectively $127$; so the search finds families when they exist.

To test that the normalisation loses no family, \texttt{brute.cpp} lists all planar difference families with three blocks of sizes $a>b>c\ge2$ in $\ZZ_p$, for every prime $p\le83$ for which such sizes exist, using only translations ($0$ in each block). For $(p;a,b,c)=(29;5,3,2)$, $(61;7,4,3)$, $(79;7,6,3)$ it finds $8$, $12$, $4$ classes, and for $(53;6,5,2)$, $(71;8,4,2)$, $(79;8,5,2)$, $(83;8,5,3)$ none. In every case \texttt{search.cpp} finds exactly the same classes (\texttt{code/brute\_vs\_search.txt}).

\section{Proof of Corollary~\ref{c:small}}

The Bruck--Ryser theorem excludes $n=6,14,21,22,30,33$. For the other non-prime-powers $n<36$, a collineation of prime order $p>n+1$ requires $p\mid n^2+n+1$ by Lemma~\ref{l:fpf}. The cases are:
\begin{center}
\begin{tabular}{rrrl}
\toprule
$n$ & $n^2+n+1$ & $p$ & excluded by \\
\midrule
10 & $3\cdot37$ & 37 & norm parity, $q=2$ \\
12 & $157$ & 157 & norm parity, $q=3$ \\
15 & $241$ & 241 & norm parity, $q=3$ \\
18 & $7^3$ & --- & no such $p$ \\
20 & $421$ & 421 & norm parity, $q=5$ \\
24 & $601$ & 601 & cyclic plane~\cite{EM} \\
26 & $19\cdot37$ & 37 & norm parity, $q=2$ \\
28 & $3\cdot271$ & 271 & Theorem~\ref{t:main} \\
34 & $3\cdot397$ & 397 & norm parity, $q=2$ \\
35 & $13\cdot97$ & 97 & norm parity, $q=5$ \\
\bottomrule
\end{tabular}
\end{center}
\emph{Norm parity}~\cite{K18}: if a collineation of odd prime order $p$ has $R$ point orbits of length $p$, and $q\ne p$ is a prime with $-1\in\langle q\bmod p\rangle$, then $R\,v_q(n)$ is even. In each marked case $v_q(n)=1$, $R=v/p$ is odd, and $-1$ is a power of $q$ modulo $p$.

For $n=24$ we have $R=1$, so $\sigma$ acts regularly on the points and the plane is cyclic; there is no cyclic projective plane of order $24$~\cite{EM}.

For $n=28$ the norm-parity lemma does not apply: $v_2(28)=2$ is even, and $7$ has order $135$ modulo $271$, so $-1\notin\langle7\bmod271\rangle$.

\section*{Files}
\texttt{code/search.cpp} (the search), \texttt{code/run.sh} (the full run), \texttt{code/results.txt} (per-chunk output), \texttt{code/classes.py} (classes of families, used for the controls), \texttt{code/recheck.cpp} and \texttt{code/stage\_counts.txt} (the re-check), \texttt{code/brute.cpp} and \texttt{code/brute\_vs\_search.txt} (the comparison with brute force).

\begin{thebibliography}{9}
\bibitem{BR} R.~H.~Bruck and H.~J.~Ryser, The nonexistence of certain finite projective planes, \emph{Canad. J. Math.} 1 (1949), 88--93.
\bibitem{EM} T.~A.~Evans and H.~B.~Mann, On simple difference sets, \emph{Sankhy\=a} 11 (1951), 357--364.
\bibitem{K18} Y.~Kobayashi, A projective plane of order 18 admits no collineation of order 17, preprint (2026), \url{https://github.com/yudaikobayashi/finite-projective-plane-order-18}, norm-parity lemma.
\end{thebibliography}

\end{document}
